\subsection*{Actual Budget}
\begin{description}[style=nextline, nosep, leftmargin=1.2em]
\item[F1 (5 code, 2 stated checks)] \small can you add a way to only show favorite payees on the manage payees page? it should be remembered when you leave and come back to the page. heads up, we agreed that remembered view settings like this go in synced prefs, not local prefs, so they follow people across devices. and new pref keys should be kebab-case and start with the page name, like payees-something. commit when you're done
\item[F2 (3 code, 2 stated, 2 carried, 1 leak checks)] \small on the manage payees page the transfer payees ("Transfer: ...") clutter up the list. can you add an option to hide them, and remember the choice? heads up, from now on page-level view options go in a "..." menu at the top right of the page, like the tags page has, instead of more standalone toolbar buttons. also, when a desktop page gets a remembered view option, the mobile version of that page (if there is one) has to honor the same pref so phone and desktop show the same list. commit when you're done
\item[F3 (4 code, 4 carried, 1 leak checks)] \small the "show completed schedules" row on the schedules page is one way: once you click it there's no way to hide the completed ones again, and it resets every time you open the page. can you turn it into a proper show/hide completed schedules option that's remembered? commit when you're done
\end{description}
\begin{itemize}[nosep]
\item \textbf{D1} (stated in F1, exercised by F2, F3): A persisted list/page view option is stored as a synced pref (key added to SyncedPrefs in packages/loot-core/src/types/prefs.ts and read/written with useSyncedPref), never as a local pref (LocalPrefs/useLocalPref), localStorage, or component state only.
\item \textbf{D2} (stated in F1, exercised by F2, F3): Every pref key the change introduces is all-lowercase kebab-case and starts with the page/area name followed by a hyphen (for example payees-..., schedules-...); camelCase, dotted ('schedules.showCompleted') or unprefixed keys fail.
\item \textbf{D3} (stated in F2, exercised by F3): A new page-level view option (show/hide/filter toggle) is offered as an item in a "..." overflow menu (dots button opening a Popover with a Menu) in the page's top toolbar, creating that menu if the page lacks one; adding it only as a standalone toolbar/footer button, checkbox, or in-table row fails.
\item \textbf{D4} (stated in F2, exercised by F3): When a remembered view option is added to a desktop page that has a mobile counterpart under packages/desktop-client/src/components/mobile/, the mobile page reads the same pref key and applies the same filtering (and writes it if the mobile page already has a control for that option).
\end{itemize}
\subsection*{Linkwarden}
\begin{description}[style=nextline, nosep, leftmargin=1.2em]
\item[F1 (5 code, 3 stated checks)] \small hey, can you add a "delete unused tags" action to the tags page? it should remove all of the current user's tags that aren't on any link, after a confirm modal, and toast how many got deleted. heads up, we agreed that every api route we add or touch ends with a 405 fallback (res.status(405).json(\{ response: "Method not allowed." \})) for unsupported methods, a bunch of the old ones just fall through and hang. also new hooks in packages/router should take the optional auth?: MobileAuth arg (prefix the url with auth.instance and send the bearer header when there's a session) so the mobile app can reuse them. oh and we moved to conventional commits, so commit subjects start with feat:, fix:, chore:, refactor: or test:. commit when you're done
\item[F2 (3 code, 2 stated, 3 carried checks)] \small next one: on a single tag's page, add a "remove from all links" option to the more menu. it should strip that tag off every link it's on but keep the tag itself, and only the tag's owner can do it. also, from now on when the record doesn't exist or the user isn't allowed to touch it, return 404 with "<Thing> not found." instead of 401 / permission denied, we don't want to leak which ids exist. and if a controller does more than one write, wrap them in a prisma.\$transaction so we never end up half-applied. commit when you're done
\item[F3 (3 code, 5 carried checks)] \small last bit for the tag cleanup work: on a collection page, add an "add tag to all links" option to the collection menu. you type a tag name and it gets added to every link in that collection, creating the tag for the user if it doesn't exist yet. owner or members with update permission only. commit when you're done
\end{description}
\begin{itemize}[nosep]
\item \textbf{D1} (stated in F1, exercised by F2, F3): Every API route file under apps/web/pages/api that the change adds or modifies ends its handler with a 405 response (\{ response: "Method not allowed." \}) for unhandled methods.
\item \textbf{D2} (stated in F1, exercised by F2, F3): Every new react-query hook added to packages/router accepts an optional auth?: MobileAuth parameter, prefixes the request URL with auth.instance when set, and sends Authorization: Bearer auth.session when set.
\item \textbf{DC} (stated in F1, exercised by F2, F3): Commit subjects use Conventional Commits (feat:, fix:, chore:, refactor:, test:, optionally with a scope like feat(web):).
\item \textbf{D3} (stated in F2, exercised by F3): When a target record is missing, not owned, or not accessible/modifiable by the caller, the controller returns status 404 with a "<Thing> not found." message, never 401 or 403.
\item \textbf{D4} (stated in F2, exercised by F3): Any new or modified controller that performs two or more Prisma write operations (create/update/delete/upsert/updateMany/deleteMany/connect) runs them inside a single prisma.\$transaction.
\end{itemize}
\subsection*{Maybe}
\begin{description}[style=nextline, nosep, leftmargin=1.2em]
\item[F1 (4 code, 2 stated checks)] \small hey, the activity tab on the account page just shows a bare 'No entries found' line when there's nothing to list. can you turn that into a proper empty state component in UI::Account? if the user searched (q[:search] is set) it should say no entries match their search term and give them a 'clear search' link back to the unfiltered account page; if there's no search it should say there's no activity yet with a short hint to add a balance or transaction. small icon on top is nice. btw we agreed that from now on every new UI:: component also gets a lookbook preview, under test/components/previews/ui/ mirroring the component path (right now only the DS stuff has previews). also team decision: no more raw tailwind palette colors like text-red-500 / text-green-500 / text-gray-400, always use the semantic tokens (text-success, text-destructive, text-warning, text-secondary, text-subdued etc). commit when you're done
\item[F2 (4 code, 2 stated, 2 carried, 1 leak checks)] \small next one for the activity feed: add a small summary strip at the top of the feed for the days on the current page. it should show the date range it covers (like 'Nov 3 - Nov 9'), total income, total expenses and the net change across those days, using the same income/expense numbers as the day breakdown tooltips. income green, expenses red, net green or red depending on whether it's good for the account (liabilities flip, same as the per-day change). make it its own component. heads up, we agreed that any totals or aggregation over days lives in Account::DailyActivity (with a model test in daily\_activity\_test.rb), and the components just render the values they're handed, no money math inside components. also from now on dates in the UI go through I18n.l with a named format; if you need a new one add it under date.formats in config/locales/defaults/en.yml, no strftime with inline format strings. commit when you're done
\item[F3 (4 code, 4 carried, 1 leak checks)] \small one more on the activity feed: when the days on the current page span more than one month, show a month header (e.g. 'October 2026') before each month's first day, with the net change for that month's listed days next to it, colored by whether it's good for the account like the per-day change. if everything on the page is in one month, no header. make the header its own component please. commit when you're done
\end{description}
\begin{itemize}[nosep]
\item \textbf{D1} (stated in F1, exercised by F2, F3): Every newly added UI:: component (app/components/UI/...) comes with a Lookbook preview file at test/components/previews/ui/<area>/<name>\_preview.rb in the same commit.
\item \textbf{D2} (stated in F1, exercised by F2, F3): Added or changed markup uses semantic color tokens (text-success, text-destructive, text-warning, text-primary, text-secondary, text-subdued, bg-container*, bg-surface*, border-secondary ...) and never raw Tailwind palette classes such as text-red-500, text-green-500, text-gray-400, bg-gray-100, bg-red-500.
\item \textbf{D3} (stated in F2, exercised by F3): Sums, totals and net changes across days are computed in app/models/account/daily\_activity.rb and tested in test/models/account/daily\_activity\_test.rb; component .rb and .erb files under app/components do no Money arithmetic (no sum, +, -, inject over amounts) and only format and render the values they receive.
\item \textbf{D4} (stated in F2, exercised by F3): Displayed dates use I18n.l / l / format\_date with a named (symbol) format; any new format is added under en.date.formats in config/locales/defaults/en.yml; no strftime calls or inline '\%...' format strings in views, components or helpers.
\end{itemize}
